\begin{abstract}
The context window of a large language model is not memory. It is
L1~cache---a small, fast, expensive resource that the field treats
as the entire memory system. There is no L2, no virtual memory, no
paging. Every tool definition, every system prompt, and every stale
tool result occupies context for the lifetime of the session. The
result is measurable: across 857~production sessions and
4.45~billion effective input tokens, 21.8\% is structural waste.

We present \textsc{Pichay}, a demand paging system for LLM context
windows. Implemented as a transparent proxy between client and
inference API, \textsc{Pichay} interposes on the message stream to
evict stale content, detect page faults when the model re-requests
evicted material, and pin working-set pages identified by fault
history. In offline replay across 1.4~million simulated evictions,
the fault rate is 0.0254\%. In live production deployment over
681~turns, the system reduces context consumption by up to 93\%
(5{,}038KB $\to$ 339KB); under extreme sustained pressure, the
system remains operational but exhibits the expected thrashing
pathology, with repeated fault-in of evicted content.

The key observation is that the problems the field faces---context
limits, attention degradation, cost scaling, lost state across
sessions---are virtual memory problems wearing different clothes.
The solutions exist: working set
theory~\citep{denning1968working}, demand paging, fault-driven
replacement policies, and memory hierarchies with multiple
eviction-managed levels. We describe the architecture of a full
memory hierarchy for LLM systems (L1~through persistent storage),
report on the first three levels deployed in production use
(L1~eviction, L2~fault-driven pinning, L3~model-initiated
conversation compaction), and identify cross-session memory as the
remaining frontier.
\end{abstract}
